\documentclass[12pt]{letter}
\usepackage[margin=1in]{geometry}
\usepackage{url}

\signature{Jeffrey D.\ Varner, Ph.D.\\
Associate Professor\\
Robert Frederick Smith School of\\
Chemical and Biomolecular Engineering\\
Cornell University\\
Ithaca, NY 14853\\
\url{jdv27@cornell.edu}}

\date{\today}

\begin{document}

\begin{letter}{Editor-in-Chief\\
Computer Methods and Programs in Biomedicine}

\opening{Dear Editor,}

I am pleased to submit the enclosed manuscript, ``ParetoEnsembles.jl: A Julia Package for Multiobjective Parameter Estimation Using Pareto Optimal Ensemble Techniques,'' for consideration as an Original Article in \textit{Computer Methods and Programs in Biomedicine}.

Mathematical models of biological systems frequently contain many adjustable parameters that must be estimated from multiple, potentially conflicting datasets. Rather than reporting a single best-fit parameter vector, generating an ensemble of near-optimal parameter sets provides richer information about what the data can and cannot constrain. This manuscript presents ParetoEnsembles.jl, an open-source Julia package that generates such ensembles using a simulated-annealing-based algorithm with several algorithmic improvements over prior implementations, including a correction from weak to strict Pareto dominance, an incremental ranking scheme that reduces per-iteration cost from quadratic to linear, and multi-chain parallel execution.

We believe this work is well suited to the scope of CMPB for three reasons. First, the package is a computational tool validated on biomedical problems: we demonstrate it on a cell-free gene expression model fitted to experimental data and on the Hockin--Mann blood coagulation cascade, a clinically relevant system with 34 species and 42 rate constants. Second, the manuscript goes beyond software description to show how parameter ensembles enable downstream biomedical analyses, including held-out prediction with uncertainty quantification, identifiability analysis through pairwise parameter correlations, and patient-specific hemophilia~A simulations. Third, we validate the approach against published experimental thrombin generation data from Butenas et al.\ (1999), demonstrating that the ensemble predicts held-out prothrombin conditions to within 1--10\% despite inherent model approximation error.

The manuscript has not been published elsewhere and is not under consideration at any other journal. A preprint is available on arXiv. All code and data needed to reproduce the results are publicly available at \url{https://github.com/varnerlab/ParetoEnsembles.jl}, and the package is registered in the Julia General registry.

Thank you for considering this submission. I look forward to your response.

\closing{Sincerely,}

\end{letter}
\end{document}
